% Passes 7138--7145: C2 normal form, matrix quotient, and semilinear hexad.
\subsection{$C_2$ normal form, finite Gram reduction, and a semilinear six-witness configuration}
\label{sec:pass7138-7145-c2-normalform}

Let $g$ be the involution stabilizing the verified $51$-point partial ovoid $S\subset W(3,9)$.  In dual Lagrangian bases $E_+\oplus E_-$ it takes the normal form
\[
 g=\operatorname{diag}(I_2,-I_2),\qquad
 \langle(u,w),(u',w')\rangle=u\cdot w'-w\cdot u'.
\]
Thus $S$ consists of one fixed point together with $25$ two-cycles.  Each nonfixed cycle lies on a unique transversal joining the two fixed generator lines.  More generally, for every odd $q$, a non-scalar projective involutory symplectic similitude of multiplier $-1$ has two two-dimensional totally isotropic eigenspaces, hence fixed locus two disjoint generator lines.  Its admissible nonfixed two-cycles number
\[
 \frac{q(q^2-1)}2,
\]
while the non-isotropic transversals number $q(q+1)$, with $(q-1)/2$ involution-pairs on each transversal.

The rank-four Gram search for a hypothetical $52$-set admits a further exact gauge quotient.  A nonsingular four-anchor block has one of only eight canonical matching-product types over $\mathbb F_9$,
\[
(1,1,2),(1,1,3),(1,1,4),(1,1,5),
(1,2,3),(1,2,4),(1,3,4),(1,3,5).
\]
After fixing a type, every remaining point is represented by a normalized row $(1,a,b,c)$ with $a,b,c\in\mathbb F_9^*$, so there are exactly $512$ row states.  Consequently a $52$-point partial ovoid exists if and only if at least one of the eight corresponding compatibility graphs contains a $48$-clique.  The known $51$-set lies in type $(1,3,5)$ and supplies an exact $47$-clique.  No $48$-clique impossibility is claimed here.

Collapsing the involution gives an all-$q$ weighted conflict graph with $2(q+1)$ fixed weight-one nodes and $q(q^2-1)/2$ admissible weight-two nodes.  Their valencies are
\[
 d_{\rm fixed}=\frac{q^2+q+2}{2},\qquad
 d_{\rm pair}=\frac{2q^2-q+3}{2}.
\]
At $q=9$ the quotient has $380$ nodes and $14500$ edges; at $q=7$ it has $184$ nodes and $4180$ edges.  A compact spectral formula has been verified computationally at $q=3,5,7,9$ but remains a conjecture rather than an all-$q$ theorem.

The transversal carrier itself has a matrix realization.  A non-isotropic endpoint pair $([u],[w])$ determines
\[
 E=\frac{w u^{\mathsf T}}{u\cdot w},
 \qquad E^2=E,\quad \operatorname{rank}E=1,\quad \operatorname{tr}E=1.
\]
Hence non-isotropic transversals are canonically the $q(q+1)$ rank-one trace-one idempotents of $M_2(\mathbb F_q)$.  This is distinct from the larger nonzero rank-one Fourier sector used in the earlier matrix-ring construction.

Finally let $F:x\mapsto x^3$ be Frobenius on $\mathbb F_9$.  The permutations induced by $g$ and $F$ generate a dihedral group of order $12$ with $|gF|=6$.  Its orbit on $S$ consists of six $51$-point partial ovoids, every pair intersecting in exactly four points.  Their union has $248$ points with multiplicities $192\cdot1+54\cdot2+2\cdot3$; the two triple points are graph indices $50$ and $80$.  The six binary incidence vectors generate a
\[
 \boxed{[248,6,51]_2}
\]
code with integer generator Gram matrix $47I_6+4J_6$ and binary Gram matrix $I_6$.  Its weight enumerator is
\[
1+6z^{51}+15z^{94}+18z^{129}+2z^{133}+9z^{156}+6z^{160}+6z^{179}+z^{194}.
\]
These are finite symplectic, matrix, and coding statements only; in particular they neither establish $\alpha(W(3,9))=51$ nor imply a physical or CSS identification.
